\section{Problem 66 --- Round 4 (continuation \& spectral-rigidity mode)}

\subsection*{1) ROUND-4 OBJECTIVE}
\textbf{Path (C): obstruction/correction direction.}  Round~3 proved that if
\[
\frac{r(n)}{\log n}\to L\in(0,\infty),\qquad r(n):=(1_A\ast 1_A)(n),
\]
then $A$ is asymptotically equidistributed in every \emph{fixed residue class}.  The exact gap left after Round~3 is that this only controls \emph{rational Fourier modes} (roots of unity).  The most promising next step is therefore to pass from rational frequencies to \emph{all fixed frequencies on the unit circle}.  

In this round I prove a strictly stronger rigidity statement: the twisted Laplace transform
\[
G_\theta(t):=\sum_{n\ge 1}1_A(n)e^{-tn}e^{in\theta}
\]
is lower-order for every fixed $\theta\notin 2\pi\mathbb Z$, and in fact this decay is \emph{uniform} away from an $o(1)$-arc around $\theta=0$.  As a consequence, the exponentially weighted measures supported on $\{n\alpha\}$ with $n\in A$ converge to Haar measure on the orbit closure of $\alpha$; in particular, for irrational $\alpha$, the set $A$ is \emph{Abel-weightedly equidistributed modulo $1$} along the rotation $n\mapsto n\alpha$.

This is a genuine extension of Round~3 from finite cyclic Fourier analysis to the full torus.

\subsection*{2) ROUND-3 FOUNDATION USED}
I use the following previously established results as black boxes.
\begin{itemize}
\item \textbf{Round~2 Lemma 66.3:}
\[
\sum_{n\ge2}(\log n)e^{-tn}\sim \frac1t\log\frac1t.
\]
\item \textbf{Round~2 Lemma 66.4:} with
\[
F(t):=\sum_{n\ge 2} r(n)e^{-tn},
\]
we have
\[
F(t)\sim \frac{L}{t}\log\frac1t\qquad(t\to 0^+).
\]
\item \textbf{Round~2 Lemma 66.5:} with
\[
G(t):=\sum_{n\ge 1}1_A(n)e^{-tn},
\]
we have
\[
G(t)\sim \sqrt{L}\,t^{-1/2}\sqrt{\log\tfrac1t}.
\]
\item \textbf{Round~3 Theorem 66.10 and Corollary 66.11:} every fixed modulus $q$ is equidistributed, equivalently every nontrivial \emph{rational} Fourier mode is asymptotically negligible in natural counting.  I use this as the structural benchmark that the present round must strictly strengthen.
\end{itemize}

\subsection*{3) NEW INSIGHT / TOOL (ROUND-4)}
\textbf{New tool: uniform summation-by-parts for twisted Laplace sums.}

For
\[
S_\theta(t):=\sum_{n\ge 2}(\log n)e^{-tn}e^{in\theta},
\]
I prove the quantitative bound
\[
S_\theta(t)\ll \frac{\log(1/t)}{\operatorname{dist}(\theta,2\pi\mathbb Z)}
\]
uniformly for $\theta\notin 2\pi\mathbb Z$.  Since the main untwisted size is
\[
\sum_{n\ge2}(\log n)e^{-tn}\asymp \frac1t\log\frac1t,
\]
this shows that twisting by any angle separated from $0$ by more than $\gg t$ forces a power-sized drop in the Laplace transform.

Inserted into the identity
\[
\sum_{n\ge2} r(n)e^{-tn}e^{in\theta}=G_\theta(t)^2,
\]
this gives a \emph{uniform spectral exclusion principle}: the only place where the generating function can have the principal singular size is inside a vanishing arc around $z=1$.

\subsection*{4) ATTACK PLAN (ROUND-4)}
The remaining gap after Round~3 is:
\begin{quote}
\emph{Round~3 controls only rational phases $e^{2\pi i m/q}$; it does not rule out a persistent irrational or quasiperiodic bias of $1_A$.}
\end{quote}

I address exactly that gap by proving the following claims.
\begin{enumerate}
\item For the model coefficients $(\log n)$, twisted Laplace sums at angle $\theta$ are only of size $O\big(\log(1/t)/\delta\big)$, where $\delta=\operatorname{dist}(\theta,2\pi\mathbb Z)$.
\item Therefore the twisted Laplace transform of the representation function,
\[
R_\theta(t):=\sum_{n\ge2}r(n)e^{-tn}e^{in\theta},
\]
is $o\big((1/t)\log(1/t)\big)$ uniformly whenever $\delta\gg t$.
\item Since $R_\theta(t)=G_\theta(t)^2$, the twisted transform $G_\theta(t)$ is $o\big(t^{-1/2}\sqrt{\log(1/t)}\big)$ uniformly away from a $t$-scale neighborhood of $0$.
\item Normalizing by the untwisted mass $G(t)$ then yields convergence of the exponentially weighted orbital measures
\[
\mu_t^{(\alpha)}:=\frac1{G(t)}\sum_{n\ge1}1_A(n)e^{-tn}\,\delta_{\{n\alpha\}}
\]
to Haar measure on the compact subgroup generated by $\alpha$.
\end{enumerate}

This strictly strengthens Round~3: rational frequencies become a special case, while irrational frequencies are genuinely new.

\subsection*{5) WORK (ROUND-4)}
Throughout, write
\[
a_n:=1_A(n),\qquad r(n):=\sum_{k=1}^{n-1}a_ka_{n-k},
\]
and assume
\begin{equation}
\label{eq:R4-assumption}
\frac{r(n)}{\log n}\to L\in(0,\infty).
\end{equation}
Also put
\[
M(t):=\frac1t\log\frac1t\qquad(t\in(0,e^{-1}]).
\]

\medskip
\noindent\textbf{Lemma 66.12 (uniform twisted model bound).}
Let
\[
S_\theta(t):=\sum_{n\ge2}(\log n)e^{-tn}e^{in\theta}
\qquad(t>0,\ \theta\in\mathbb R),
\]
and let
\[
\delta(\theta):=\operatorname{dist}(\theta,2\pi\mathbb Z)\in[0,\pi].
\]
Then for every $t\in(0,e^{-1}]$ and every $\theta\notin 2\pi\mathbb Z$,
\[
S_\theta(t)\ll \frac{\log(1/t)}{\delta(\theta)},
\]
with an absolute implied constant.

\begin{proof}
Fix $t$ and $\theta\notin 2\pi\mathbb Z$, and set
\[
g_t(n):=(\log n)e^{-tn}.
\]
Let
\[
B_\theta(N):=\sum_{n=2}^{N} e^{in\theta}.
\]
Since this is a finite geometric sum,
\[
|B_\theta(N)|\le \frac{2}{|1-e^{i\theta}|}
=\frac{1}{|\sin(\theta/2)|}
\ll \frac{1}{\delta(\theta)},
\]
because $|\sin u|\ge 2|u|/\pi$ for $|u|\le \pi/2$ and $\delta(\theta)\in[0,\pi]$.

Now summation by parts gives, for every $N\ge2$,
\[
\sum_{n=2}^{N} e^{in\theta}g_t(n)
=B_\theta(N)g_t(N)-\sum_{n=2}^{N-1} B_\theta(n)\big(g_t(n+1)-g_t(n)\big).
\]
Since $g_t(N)\to0$ as $N\to\infty$, letting $N\to\infty$ yields
\[
|S_\theta(t)|
\ll \frac{1}{\delta(\theta)}\sum_{n\ge2}|g_t(n+1)-g_t(n)|.
\]
By the fundamental theorem of calculus,
\[
|g_t(n+1)-g_t(n)|\le \int_n^{n+1}|g_t'(x)|\,dx,
\]
so
\[
\sum_{n\ge2}|g_t(n+1)-g_t(n)|\le \int_1^{\infty}|g_t'(x)|\,dx.
\]
As in Round~2 Lemma~66.3,
\[
g_t'(x)=\Big(\frac1x-t\log x\Big)e^{-tx},
\]
and therefore
\[
\int_1^{\infty}|g_t'(x)|\,dx
\le \int_1^{\infty}\Big(\frac1x+t\log x\Big)e^{-tx}\,dx
\ll \log\frac1t.
\]
Combining the estimates gives
\[
S_\theta(t)\ll \frac{\log(1/t)}{\delta(\theta)},
\]
as required.
\end{proof}

\medskip
\noindent\textbf{Proposition 66.13 (uniform twisted cancellation for the representation transform).}
Define
\[
R_\theta(t):=\sum_{n\ge2} r(n)e^{-tn}e^{in\theta}
\qquad(t>0,\ \theta\in\mathbb R).
\]
Let $\eta:(0,e^{-1}]\to(0,\pi]$ satisfy
\[
\frac{\eta(t)}{t}\to\infty\qquad(t\to0^+).
\]
Then
\[
\sup_{\delta(\theta)\ge \eta(t)} \frac{|R_\theta(t)|}{M(t)}\to 0
\qquad(t\to0^+).
\]
In particular, for every fixed $\theta\notin2\pi\mathbb Z$,
\[
R_\theta(t)=o\!\left(\frac1t\log\frac1t\right).
\]

\begin{proof}
Fix $\varepsilon\in(0,1)$.  By \eqref{eq:R4-assumption}, there exists $N_\varepsilon$ such that for all $n\ge N_\varepsilon$,
\[
|r(n)-L\log n|\le \varepsilon\log n.
\]
Split
\[
R_\theta(t)
=\sum_{2\le n< N_\varepsilon} r(n)e^{-tn}e^{in\theta}
+L\sum_{n\ge N_\varepsilon}(\log n)e^{-tn}e^{in\theta}
+\sum_{n\ge N_\varepsilon}\big(r(n)-L\log n\big)e^{-tn}e^{in\theta}.
\]
Hence, uniformly in $\theta$,
\[
|R_\theta(t)|
\le O_\varepsilon(1)
+L\left|\sum_{n\ge N_\varepsilon}(\log n)e^{-tn}e^{in\theta}\right|
+\varepsilon\sum_{n\ge N_\varepsilon}(\log n)e^{-tn}.
\]
By Lemma~66.12, deleting finitely many terms changes the twisted sum by $O_\varepsilon(1)$, so
\[
\sup_{\delta(\theta)\ge \eta(t)}
\left|\sum_{n\ge N_\varepsilon}(\log n)e^{-tn}e^{in\theta}\right|
\ll \frac{\log(1/t)}{\eta(t)}.
\]
Also, by Round~2 Lemma~66.3,
\[
\sum_{n\ge N_\varepsilon}(\log n)e^{-tn}
\sim \frac1t\log\frac1t = M(t).
\]
Therefore
\[
\sup_{\delta(\theta)\ge \eta(t)} |R_\theta(t)|
\le O_\varepsilon(1)+O\!\left(\frac{\log(1/t)}{\eta(t)}\right)+\varepsilon(1+o(1))M(t).
\]
Divide by $M(t)=(1/t)\log(1/t)$:
\[
\sup_{\delta(\theta)\ge \eta(t)} \frac{|R_\theta(t)|}{M(t)}
\le o(1)+O\!\left(\frac{t}{\eta(t)}\right)+\varepsilon(1+o(1)).
\]
Since $\eta(t)/t\to\infty$, the middle term tends to $0$.  Thus
\[
\limsup_{t\to0^+}
\sup_{\delta(\theta)\ge \eta(t)} \frac{|R_\theta(t)|}{M(t)}
\le \varepsilon.
\]
Because $\varepsilon>0$ is arbitrary, the stated uniform $o(M(t))$ follows.

For fixed $\theta\notin 2\pi\mathbb Z$, choose for instance $\eta(t)=\tfrac12\delta(\theta)$ for all sufficiently small $t$; then $\eta(t)/t\to\infty$, so the pointwise statement follows.
\end{proof}

\medskip
\noindent\textbf{Theorem 66.14 (uniform spectral exclusion away from $1$).}
Let
\[
G_\theta(t):=\sum_{n\ge1} a_n e^{-tn}e^{in\theta}
\qquad(t>0,\ \theta\in\mathbb R).
\]
Then
\[
R_\theta(t)=G_\theta(t)^2
\qquad\text{for all }t>0,\ \theta\in\mathbb R.
\]
Moreover, for every function $\eta(t)$ with $\eta(t)/t\to\infty$,
\[
\sup_{\delta(\theta)\ge \eta(t)}
\frac{|G_\theta(t)|}{t^{-1/2}\sqrt{\log(1/t)}}\to0
\qquad(t\to0^+).
\]
In particular, for every fixed $\theta\notin 2\pi\mathbb Z$,
\[
G_\theta(t)=o\!\left(t^{-1/2}\sqrt{\log\tfrac1t}\right).
\]

\begin{proof}
The convolution identity is immediate from absolute convergence:
\[
G_\theta(t)^2
=\sum_{k,\ell\ge1} a_ka_\ell e^{-t(k+\ell)}e^{i(k+\ell)\theta}
=\sum_{n\ge2}\Big(\sum_{k=1}^{n-1}a_ka_{n-k}\Big)e^{-tn}e^{in\theta}
=R_\theta(t).
\]
Hence
\[
|G_\theta(t)|^2=|R_\theta(t)|.
\]
By Proposition~66.13,
\[
\sup_{\delta(\theta)\ge \eta(t)} |R_\theta(t)| = o(M(t)).
\]
Taking square roots gives
\[
\sup_{\delta(\theta)\ge \eta(t)} |G_\theta(t)| = o\big(M(t)^{1/2}\big)
=o\!\left(t^{-1/2}\sqrt{\log\tfrac1t}\right),
\]
which is exactly the claim.
\end{proof}

\medskip
\noindent\textbf{Corollary 66.15 (Abel-weighted orbit equidistribution on the circle).}
For $\alpha\in\mathbb R$, define a probability measure on $\mathbb T:=\mathbb R/\mathbb Z$ by
\[
\mu_t^{(\alpha)}:=\frac{1}{G(t)}\sum_{n\ge1} a_n e^{-tn}\,\delta_{\{n\alpha\}}
\qquad(t>0),
\]
where $\{x\}$ denotes the class of $x$ mod $1$.  Let
\[
H_\alpha:=\overline{\{n\alpha\pmod 1:n\in\mathbb Z\}}\le \mathbb T
\]
be the closed subgroup generated by $\alpha$, and let $m_{H_\alpha}$ denote Haar probability measure on $H_\alpha$.

Then
\[
\mu_t^{(\alpha)}\Longrightarrow m_{H_\alpha}
\qquad(t\to0^+)
\]
in the weak-* topology.  Equivalently, for every continuous $1$-periodic function $f$,
\[
\frac{\sum_{n\ge1} a_n e^{-tn}f(n\alpha)}{\sum_{n\ge1} a_n e^{-tn}}
\to \int_{H_\alpha} f\,dm_{H_\alpha}.
\]
In particular:
\begin{itemize}
\item if $\alpha\notin\mathbb Q$, then $H_\alpha=\mathbb T$ and
\[
\frac{\sum_{n\ge1} a_n e^{-tn}f(n\alpha)}{\sum_{n\ge1} a_n e^{-tn}}
\to \int_0^1 f(x)\,dx;
\]
\item if $\alpha=p/q\in\mathbb Q$ in lowest terms, then $H_\alpha=\{0,1/q,\dots,(q-1)/q\}$ and the limit measure is uniform on this $q$-point orbit.
\end{itemize}

\begin{proof}
For each integer $m$, the $m$-th Fourier coefficient of $\mu_t^{(\alpha)}$ is
\[
\widehat\mu_t^{(\alpha)}(m)
:=\int_{\mathbb T} e^{2\pi i mx}\,d\mu_t^{(\alpha)}(x)
=\frac{1}{G(t)}\sum_{n\ge1} a_n e^{-tn}e^{2\pi i mn\alpha}
=\frac{G_{2\pi m\alpha}(t)}{G(t)}.
\]
If $m\alpha\in\mathbb Z$, then $G_{2\pi m\alpha}(t)=G(t)$, so
\[
\widehat\mu_t^{(\alpha)}(m)=1.
\]
If $m\alpha\notin\mathbb Z$, then $2\pi m\alpha\notin2\pi\mathbb Z$, and Theorem~66.14 gives
\[
\widehat\mu_t^{(\alpha)}(m)=\frac{G_{2\pi m\alpha}(t)}{G(t)}\to0,
\]
because Round~2 Lemma~66.5 gives $G(t)\sim \sqrt{L}\,t^{-1/2}\sqrt{\log(1/t)}$.

Thus, for every $m\in\mathbb Z$,
\[
\widehat\mu_t^{(\alpha)}(m)\to
\begin{cases}
1,& m\in H_\alpha^{\perp},\\
0,& m\notin H_\alpha^{\perp},
\end{cases}
\]
where $H_\alpha^{\perp}=\{m\in\mathbb Z: e^{2\pi i mx}=1\ \forall x\in H_\alpha\}$ is the annihilator of $H_\alpha$.  These are exactly the Fourier coefficients of Haar probability measure on $H_\alpha$.  Since trigonometric polynomials are dense in $C(\mathbb T)$, the Fourier coefficients determine the weak-* limit, proving
\[
\mu_t^{(\alpha)}\Longrightarrow m_{H_\alpha}.
\]
The stated formulation for continuous $f$ is exactly the definition of weak-* convergence.  The irrational and rational cases are the standard descriptions of $H_\alpha$.
\end{proof}

\subsection*{6) ADVERSARIAL VERIFICATION}
\begin{itemize}
\item \textbf{Interaction with Round~3.}  Round~3 handled rational frequencies by discrete Fourier analysis modulo $q$.  The present round does \emph{not} re-prove that theorem; it extends the cancellation mechanism to all real frequencies via a different summation-by-parts argument.  Rational phases become a special case of Theorem~66.14.
\item \textbf{Uniformity near $\theta=0$.}  The bound in Lemma~66.12 necessarily degenerates like $1/\delta(\theta)$.  Accordingly, Theorem~66.14 only claims uniform decay when $\delta(\theta)\gg t$; no false uniformity across a $t$-scale arc around $0$ is asserted.
\item \textbf{Ordered convolution.}  The identity $R_\theta(t)=G_\theta(t)^2$ depends on the ordered representation function.  For unordered representations the transform would differ by a diagonal correction, so the present constants and formulas are specific to the ordered problem, consistently with earlier rounds.
\item \textbf{Complex square-root issue.}  No branch choice is used.  We only infer
\[
|G_\theta(t)|^2=|R_\theta(t)|,
\]
so the uniform $o$-estimate follows from absolute values alone.
\item \textbf{Haar-limit corollary.}  The proof uses only Fourier coefficients of probability measures on the compact group $\mathbb T$ and density of trigonometric polynomials in $C(\mathbb T)$.  No unjustified de-Abelization is claimed.
\item \textbf{What is and is not proved.}  Corollary~66.15 is an \emph{Abel-weighted} equidistribution theorem.  It does \emph{not} imply that
\[
\frac{1}{A(x)}\sum_{\substack{n\le x\\ n\in A}} e^{2\pi i mn\alpha}\to0
\]
for irrational $\alpha$; such a Tauberian de-Abelization step remains open here.
\item \textbf{Consistency with earlier rounds.}  If $\alpha=p/q$, Corollary~66.15 gives only exponential-weight equidistribution on the $q$-point orbit.  Round~3 Theorem~66.10 is stronger in that special case, because it upgrades the conclusion to ordinary counting in each residue class mod~$q$.
\end{itemize}

\subsection*{7) FINAL (EXACTLY ONE)}
\textbf{UNRESOLVED (BUT STRICTLY ADVANCED).}

Round~4 proves a new rigidity theorem beyond Round~3:
\[
\sum_{n\ge1}1_A(n)e^{-tn}e^{in\theta}
=o\!\left(t^{-1/2}\sqrt{\log\tfrac1t}\right)
\qquad(\theta\notin2\pi\mathbb Z),
\]
and in fact the decay is uniform outside every arc of width $\eta(t)$ with $\eta(t)\gg t$.  Equivalently, the principal singular behavior of the generating function is spectrally localized at $z=1$; every other fixed boundary phase is negligible.

As a consequence, for every fixed $\alpha$, the exponentially weighted orbital measures supported on $\{n\alpha:n\in A\}$ converge to Haar measure on the subgroup generated by $\alpha$.  In particular, every irrational rotation is uniformly distributed along $A$ in Abel means.  This rules out an entire class of constructions with persistent fixed-frequency or quasiperiodic bias.

The original existence/nonexistence problem remains open: I have not produced either a contradiction or a construction of such a set $A$.

\subsection*{8) COMPLETION ESTIMATE (MANDATORY)}
\noindent\textbf{COMPLETION: 63\%}

\subsection*{9) REFERENCES}
\begin{itemize}
\item N. H. Bingham, C. M. Goldie, J. L. Teugels, \emph{Regular Variation}, Cambridge Univ. Press.
\item W. Rudin, \emph{Fourier Analysis on Groups}, Wiley.
\end{itemize}
